\section{Summary}\label{ht:sec:summary}

\paragraph{What is computed.} The long-range Ising (LRI) fixed point in $d$ dimensions is reached from the generalized
free field (GFF) $\phi$ of dimension $\Ds$ by the relevant $\phi^4$
deformation~\cite{ht:Fisher:1972zz,ht:Sak:1973oqx,ht:Paulos:2015jfa,ht:Behan:2017emf}. In radial quantization the GFF
becomes a set of free oscillators. We diagonalize the deformed Hamiltonian in a truncated Fock space and tune its two
couplings, the mass $r$ and the quartic $g$, so that two levels take their conformal values: $\sigma=\phi$ at $\Ds$
and its shadow at $d-\Ds$. The first excited $\mathbb Z_2$-even scalar level is then the prediction for the
dimension $\DE$ of the energy operator $\varepsilon$ (Section~\ref{ht:sec:model}). The method is Hamiltonian truncation
in radial quantization~\cite{ht:Hogervorst:2014rta,ht:Rychkov:2014eea}, applied to a nonlocal GFF.

\paragraph{The three ingredients.}
\begin{enumerate}
\item \emph{Zero mode and bath} (Section~\ref{ht:sec:split}). The zero mode costs no level, so its quanta are not capped
 by the level cutoff, and at strong coupling its fluctuations are large. The full Hamiltonian, including the
 couplings of the zero mode to the bath, is diagonalized on the product of $n_z$ zero-mode states per parity, each a
 superposition of many zero-mode quanta, with the bath states of total level $\ell\le N$. In $d=1$ this zero-mode
 basis reproduces the $\DE$ of a converged zero mode to $\HTq{S9}$ (Section~\ref{ht:sec:zerobasis}); only the bath
 cutoff $N$ is extrapolated.
\item \emph{Window correction} (Section~\ref{ht:sec:window}). A bath state of level $\ell$ can host vacuum fluctuations
 only up to level $N-\ell$. The raw spectrum therefore violates the exact descendant spacings by amounts that do not
 decay with $N$. We add to every bath state the difference $D(\ell)=E_v(N)-E_v(N-\ell)$, where $E_v(M)$ is the
 vacuum energy at bath cutoff $M$ of the corrected theory itself, zero mode included. With an energy cutoff, the
 same missing states were added back order by order with perturbative counterterms in $\phi^4$ theory in $2+1$
 dimensions~\cite{ht:Elias-Miro:2020qwz}.
\item \emph{Zero-mode basis} (Section~\ref{ht:sec:zerobasis}). If the $n_z$ states are the lowest eigenstates of the
 zero-mode part $H_z$ of $H$, the number needed at strong coupling grows with $N$, and the zero mode then sets the
 cost. We keep instead the leading eigenvectors of a reduced density matrix of the states the calculation needs,
 trained once per $\Ds$.
\end{enumerate}

\paragraph{Results.}
\begin{itemize}
\item In $d=1$ (Section~\ref{ht:sec:d1}, Table~\ref{ht:tab:d1}) the cutoff reaches $N=\HTq{S2}$ at all 13 values
 $0.001\le\Ds\le0.249$, and the series are extrapolated with a fixed family of fits.
 \begin{itemize}
 \item At $\Ds=0.249$ the result differs from the $\eps$ expansion by $\HTq{P1}$, well below the $\HTq{P4}$ to
 which we trust this check.
 \item The values are \emph{final} for $\Ds\geq0.1$ as limits of the window-corrected model. Their half-spreads
 are at most $\HTq{D14}$; at $\Ds\ge0.15$ we judge the total error to be about $\HTq{D12}$.
 \item The quoted errors are half-spreads of the extrapolations in the cutoff, not standard deviations, and they do
 not include the dependence on the window prescription. That dependence is $\HTq{M10}$ at $0.075\le\Ds\le0.125$,
 as large as the gap there to the numerical renormalization group (NRG) of the sub-ohmic spin-boson
 model~\cite{ht:Bulla:2003rns,ht:Bulla:2005bnrg} (Section~\ref{ht:sec:d1-model}).
 \item At $\Ds=0.075$ the error is \emph{one-sided}, and at $0.05$ the extrapolation is \emph{unstable}.
 \item At $\Ds\le0.025$ the corrected sequence turns over and keeps falling, so it has \emph{no limit}. This is a
 property of the window-corrected model, not of the numerics.
 \end{itemize}
\item In $d=2$ (Section~\ref{ht:d2:sec}, $N\le24$) and $d=3$ (Section~\ref{ht:d3:sec}, $N\le20$) the cutoffs are too few
 for a controlled extrapolation, and we report finite-$N$ behaviour with status words.
\end{itemize}

\paragraph{Checks.}
\begin{itemize}
\item Separately written solvers, built from the definition of the model without using the production code, reproduce
 the production spectra in each dimension (Section~\ref{ht:sec:checks}). The $d=1$ and $d=3$ solvers were written blind
 to the production code. The $d=2$ solver is not blind: its author had earlier prepared the production code for the
 cluster.
\item These checks verify that we compute the stated model. Whether its $N\to\infty$ limit is the CFT is a separate
 question, and the answer is negative in parts of the range (Section~\ref{ht:sec:caveats}).
\end{itemize}

\paragraph{Status words.} \emph{Final}, \emph{one-sided}, \emph{unstable}, \emph{no limit} and \emph{finite-$N$}
are part of the results (Section~\ref{ht:sec:d1-extrap}). A quoted error is a half-spread: half the difference between
the largest and the smallest limit of the accepted fits. It is a model spread, not a standard deviation, and it
covers the extrapolation only.
